\section{推进系统建模}
\label{sec:propulsion}

推进系统的建模是整机动力学分析的起点。本节从最基本的流体力学原理出发，
建立推力与反扭矩的数学描述，并分析电机动态响应和旋转部件的陀螺效应。

\subsection{桨盘动量理论}

动量理论（momentum theory）或称为Froude-Rankine理论，是分析螺旋桨/旋翼性能的最基本工具。
该理论将螺旋桨理想化为无限薄桨盘（actuator disk），忽略桨叶的具体几何形状，
将问题简化为控制体上的质量、动量和能量守恒\cite{leishman2006helicopter}。

\begin{figure}[H]
\centering
\resizebox{0.88\linewidth}{!}{%
\begin{tikzpicture}[>=Stealth,
  stream/.style={draw=black!55, line width=0.8pt},
  streamline/.style={draw=black!35, densely dashed, line width=0.55pt},
  section/.style={draw=black!65, line width=0.75pt},
  velocity/.style={->, draw=blue!70!black, line width=1.05pt, >={Stealth[length=2.8mm]}},
  pressure/.style={draw=red!70!black, line width=1.05pt},
  annot/.style={font=\footnotesize\sffamily, fill=white, inner sep=1.2pt},
]
  % 流管：速度沿程增加，因此截面积沿程减小
  \draw[stream] (-0.2,1.35) .. controls (1.7,1.28) and (2.6,1.02) .. (3.2,0.9)
                .. controls (4.4,0.67) and (5.6,0.57) .. (6.8,0.55);
  \draw[stream] (-0.2,-1.35) .. controls (1.7,-1.28) and (2.6,-1.02) .. (3.2,-0.9)
                .. controls (4.4,-0.67) and (5.6,-0.57) .. (6.8,-0.55);
  \draw[streamline] (-0.2,0.68) .. controls (1.7,0.64) and (2.6,0.51) .. (3.2,0.45)
                     .. controls (4.4,0.34) and (5.6,0.29) .. (6.8,0.28);
  \draw[streamline] (-0.2,-0.68) .. controls (1.7,-0.64) and (2.6,-0.51) .. (3.2,-0.45)
                     .. controls (4.4,-0.34) and (5.6,-0.29) .. (6.8,-0.28);

  % 三个截面
  \draw[section] (0,1.35) -- (0,-1.35);
  \draw[section, line width=1.5pt, fill=black!10, fill opacity=.45] (3.2,0.9) -- (3.2,-0.9);
  \draw[section] (6.6,0.56) -- (6.6,-0.56);
  \node[annot, above=2pt, align=center] at (0,1.35) {截面0\\$A_0$};
  \node[annot, above=2pt, align=center] at (3.2,0.9) {桨盘截面1\\$A$};
  \node[annot, above=2pt, align=center] at (6.6,0.56) {截面2\\$A_2$};

  % 速度箭头，长度随速度增加
  \draw[velocity] (-0.85,0) -- (0.45,0) node[annot, midway, above=2pt] {$v_0$};
  \draw[velocity] (2.15,0) -- (3.75,0) node[annot, midway, above=2pt] {$v_p=v_0+v_i$};
  \draw[velocity] (4.65,0) -- (6.75,0) node[annot, midway, above=2pt] {$v_2=v_0+2v_i$};
  \node[font=\scriptsize\sffamily, text=black!60, fill=white, inner sep=1pt] at (3.2,-1.16) {无限薄桨盘};

  % 独立压强示意，避免与流管重叠
  \draw[->, draw=black!55, line width=.55pt] (-0.1,-2.35) -- (7.05,-2.35) node[right, font=\scriptsize\sffamily] {$x$};
  \draw[->, draw=black!55, line width=.55pt] (-0.1,-2.35) -- (-0.1,-1.30)
    node[above, font=\scriptsize\sffamily] {$p-p_\infty$};
  \draw[pressure, dashed] (-0.1,-1.88) -- (7.0,-1.88);
  \node[font=\scriptsize\sffamily, text=red!70!black, fill=white, inner sep=1pt] at (0.45,-1.73) {$p_\infty$};
  \draw[pressure] (0,-1.88) .. controls (1.5,-1.9) and (2.55,-2.12) .. (3.18,-2.16);
  \draw[pressure] (3.2,-1.42) .. controls (4.0,-1.48) and (5.35,-1.75) .. (6.6,-1.88);
  \draw[pressure, ->] (3.2,-2.16) -- (3.2,-1.42);
  \node[annot, text=red!70!black, anchor=west] at (3.38,-1.7) {$\Delta p=T/A$};
\end{tikzpicture}%
}
\caption{理想桨盘的流管与压强分布示意。由于$v_0<v_p<v_2$，不可压定常流的流管沿流向收缩（$A_0>A>A_2$）；速度穿过无限薄桨盘时连续，而静压在盘面发生跃变$\Delta p=T/A$，并在远尾流中恢复至环境压强。图中 $x$ 为桨盘局部流向轴，不代表构型主图的相机左右；几何尺寸与箭头长度仅作定性示意。}
\label{fig:actuator_disk}
\end{figure}


取桨盘面积$A = \pi D^2 / 4$（$D$为桨径），定义远前方截面0（来流条件）、
桨盘截面和远后方尾流截面2构成流管控制体。假设流体无黏、不可压，
盘上压强跳变但速度连续（流量守恒要求截面相等则速度必连续），
盘前后压差在控制体轴向的积分为推力。

由质量守恒（流过各截面的质量流量相等）和轴向动量定理，
取控制体包含盘前和盘后全部流体：
\begin{equation}
    \dot{m} = \rho A v_p, \quad
    T = \dot{m} (v_2 - v_0)
    \label{eq:momentum_basic}
\end{equation}
其中$v_p$为盘面处流速，$v_0$为远前方来流速度（对于前飞机体即前行速度），
$v_2$为远后方尾流速度。

在盘面两侧分别应用伯努利方程（盘前：来流到盘面前；盘后：盘面后到尾流），
注意到盘面处速度连续但压强跳变$\Delta p = T/A$，可导出尾流速度为来流速度加两倍诱导速度：
$v_2 = v_0 + 2v_i$，其中诱导速度$v_i = v_p - v_0$为盘面处因桨盘做功而额外加速的流速。
代入式(\ref{eq:momentum_basic})即得推力的通用表达式：
\begin{equation}
    T = 2\rho A (v_0 + v_i) v_i
    \label{eq:thrust_general}
\end{equation}

这一表达式揭示了螺旋桨推力的基本物理：推力正比于空气质量流量（$\rho A v_p$）与速度增量（$v_2-v_0$）的乘积，
而速度增量恰为$2v_i$。增大推力有两种途径——增大桨盘面积（更高效地加速更多空气）或增大诱导速度（更快速地将空气向后抛）。
前者通过增大桨径实现，后者通过提高转速或增大桨距实现。

\textbf{悬停特例}（$v_0 = 0$）是分析推进系统基本性能的重要基准：
\begin{equation}
    T = 2\rho A v_i^2 \quad\Longrightarrow\quad
    v_i = \sqrt{\frac{T}{2\rho A}}
    \label{eq:hover_induced}
\end{equation}
悬停时的理想诱导功率$P_i = T v_i = T^{3/2} / \sqrt{2\rho A}$。
由此可得一个关键推论：在给定推力下，诱导功率与桨盘面积的关系为$P_i \propto 1/\sqrt{A} \propto 1/D$——
增大桨径是降低悬停功耗最有效的途径。这就是直升机旋翼大直径和多旋翼多桨盘设计的底层物理动机。
同时，理想悬停效率（figure of merit，$FM = P_i / P_{\text{actual}}$）是衡量
真实螺旋桨偏离理想动量理论程度的重要指标，在选型分析中应予以关注。

\textbf{前飞修正}（$v_0 > 0$）：在$v_0 \gg v_i$的前飞巡航工况，
诱导速度随前飞速度单调下降，推力公式趋近于$T \approx 2\rho A v_0 v_i$。
这意味着在同样的推力和桨盘面积下，前飞所需的诱导速度远小于悬停——换言之，
螺旋桨在前飞时比悬停时效率更高。这也解释了为什么固定翼飞机的巡航推进效率
远高于旋翼飞行器的悬停效率，以及为什么倾转旋翼构型在巡航模式下提供净收益。

对于拉力式（tractor）布局，前飞还带来一个额外效应：螺旋桨滑流加速流经机翼和机身的局部流速，
改变了当地动压分布，可能在一定程度上影响升力和力矩。在初步概念设计中，
滑流效应通常作为附加修正项处理，不在零阶动力学模型中展开。

\subsection{叶素理论：$k_T$、$k_Q$的物理来源}

动量理论将螺旋桨抽象为零厚度的压强跳变盘，不涉及桨叶的具体几何（弦长分布、桨距角、
翼型特性等）。叶素理论（blade element theory）\cite{mccormick1995aerodynamics}
弥补了这一不足：将桨叶分解为展向微元$dr$，每个微元视为以当地速度运动的二维翼型段，
分别计算其升力和阻力，然后沿展向积分得到总推力$T$和总扭矩$Q$。

对展向位置$r$处的微元，当地来流是旋转线速度$\omega r$（切向）与轴向速度
$v_0+v_i$（轴向，含诱导速度修正）的合成。微元产生的升力$dL$和阻力$dD$
按当地迎角（几何桨距角减流入角）由翼型气动特性给出。
将$dL$和$dD$分别投影到推力方向（轴向）和旋转平面（切向），沿桨叶展向积分，
再乘以桨叶数，得总推力和总扭矩。

将积分结果作量纲化处理（除以$\rho A (\omega R)^2$和$\rho A (\omega R)^2 R$等参考量），
引入无量纲推力系数$C_T$和扭矩系数$C_Q$：
\begin{equation}
    T = \frac{C_T \rho D^4}{4\pi^2} \omega^2 \equiv k_T \omega^2, \qquad
    Q = \frac{C_Q \rho D^5}{4\pi^2} \omega^2 \equiv k_Q \omega^2
    \label{eq:similarity}
\end{equation}
其中$C_T$、$C_Q$为无量纲系数，$D$为桨径，$\omega$为角速度。

式(\ref{eq:similarity})的形式揭示了螺旋桨性能的相似律。在固定桨距和固定前进比下，
$C_T$和$C_Q$近似为常数，推力与转速的平方成正比（$T \propto \omega^2$），
扭矩亦与转速的平方成正比（$Q \propto \omega^2$）。当桨距可调时，$C_T$和$C_Q$
均随桨距角近似线性增长——这意味着单位转速的推力和扭矩均可通过桨距调节来改变。

从量纲看：$k_T \propto \rho D^4$（推力系数$\times$密度$\times$特征面积），
$k_Q \propto \rho D^5$（扭矩系数$\times$密度$\times$特征面积$\times$特征长度）。
因此$k_Q / k_T \propto D$——反扭矩与推力之比正比于桨径。
这一量级关系对于评估滚转控制权限至关重要：差速反扭的滚转力矩量级为$k_Q \Delta(\omega^2) \sim k_Q \omega_0^2 \Delta\omega$，
而推力矢量力臂的俯仰/偏航力矩量级为$k_T \omega_0^2 \cdot a \cdot \delta$（$a$为力臂）。
两者之比约为$(k_Q / k_T) / a \sim D / a$，即桨径与力臂之比。

需要强调的是，$C_T$和$C_Q$并非真正的常数——它们随前进比$J = v_0/(nD)$（来流速度与桨尖速度之比）
和桨距角变化。在大包线（悬停$\to$高速巡航）过渡中，$C_T(J)$和$C_Q(J)$的变化可达$20\%$--$40\%$量级。
集总参数模型中的$k_T$和$k_Q$实质上是工作点处的局部线性化值，
这是该模型的主要近似来源之一。在概念级仿真中将$k_T$和$k_Q$取为固定的集总常数是工程上可接受的简化；
若需更高精度，可引入$J$相关的修正函数$k_T(J)$和$k_Q(J)$。

\subsection{反扭矩的物理机制与角动量守恒}

反扭矩（reaction torque）是螺旋桨驱动系统施加于机体的等大反向力矩。
深入理解反扭矩的物理来源对于设计差速滚转控制系统是必要的。

取转子和定子分别为研究对象的分析方法更为透彻。转子包括电机转子（永磁体）和螺旋桨组件，
对自身转轴的转动惯量为$J_p$，角速度为$\omega$。定子即固连于机体的电机定子（线圈）。
转子承受的外力矩有两个来源：定子通过电磁场施加的驱动力矩$\tau_e$（内部作用力，$+\tau_e$作用于转子），
以及空气作用于桨叶的气动阻力矩$Q$（外部力，$-Q$作用于转子，因为阻力矩方向与旋转方向相反）\cite{ducard2011fault}。

对转子列写角动量定理（转子自身转轴方向）：
\begin{equation}
    J_p \dot{\omega} = \tau_e - Q \quad\Longrightarrow\quad
    \tau_e = Q + J_p \dot{\omega}
    \label{eq:rotor_angular}
\end{equation}
其物理含义是：电磁驱动力矩必须同时克服气动阻力矩和转子加速所需的惯性力矩。

由牛顿第三定律，定子承受等大反向的力矩$-\tau_e$。从机体的视角看，
\textbf{机体承受的反扭矩方向与转子旋向相反}，大小为：
\begin{equation}
    \boxed{\tau_{\text{机体}} = \tau_e = k_Q \omega^2 + J_p \dot{\omega}}
    \label{eq:reaction_torque}
\end{equation}

此式揭示了反扭矩的两个组成部分及其物理性质：
\begin{itemize}
    \item \textbf{稳态反扭矩}$k_Q \omega^2$：在任何恒定转速下均存在，是螺旋桨将角动量持续传递给尾流所产生的反作用。
          方向始终与旋向相反，大小正比于$\omega^2$——转速加倍，反扭矩增至四倍；
    \item \textbf{惯性反扭矩}$J_p \dot{\omega}$：仅在转速变化时出现。加速阶段（$\dot{\omega} > 0$），
         惯性反扭矩叠加在稳态分量之上，机体瞬时承受更大的反扭；减速阶段（$\dot{\omega} < 0$），
         惯性反扭矩部分抵消气动阻力矩，机体承受的净反扭减小。
\end{itemize}

惯性反扭矩对飞行品质的影响不容忽视。在快速机动过程中，差速指令导致的转速剧烈变化
会通过惯性反扭矩在滚转通道产生瞬态扰动——因为差速本质上就是一台电机加速、另一台减速，
两者的$J_p \dot{\omega}$项叠加（加速电机惯性反扭增大，减速电机反扭减小），
在滚转轴上产生双倍的惯性反扭扰动。因此，\textbf{差速通道的带宽受限于电机时间常数和转子惯量}，
过快的差速变化不仅无法有效产生滚转力矩（电机来不及响应），还会引入额外的惯性扰动。

从更本质的角动量守恒视角看：将转子+排出流体+机体视为一个系统，三者角动量的变化率之和
在任何时刻精确为零。转子加速意味着转子角动量增大，同时排出流体的角动量流率$Q$也相应增大，
机体角动量（滚转速度）相应变化以维持总角动量守恒。这一视角有助于理解机动过程中的
角动量流动路径，特别是在分析快速滚转机动时的能量分配。

\subsection{电机一阶动态模型}

直流无刷电机（BLDC）的完整动态模型涉及电磁场、换相逻辑和功率电子器件的复杂交互。
对于飞行力学仿真和控制律设计，可以采用一阶惯性环节近似\cite{ducard2011fault}：
\begin{equation}
    \tau_m \dot{\omega} + \omega = \omega_{\text{cmd}}
    \label{eq:motor_lag}
\end{equation}
其中$\tau_m$为电机综合时间常数（电机+桨负载），$\omega_{\text{cmd}}$为目标转速指令。

$\tau_m$的物理构成由两部分决定：电机的电磁阻尼$k_t k_e / R$（$k_t$为扭矩常数，
$k_e$为反电动势常数，$R$为绕组电阻）提供基础的速度反馈；
螺旋桨的气动负载斜率$2 k_Q \omega_0$（随工作点增大）附加了额外的等效阻尼。
综合时间常数为：
\begin{equation}
    \tau_m = \frac{J_p}{k_t k_e / R + 2 k_Q \omega_0}
    \label{eq:tau_m}
\end{equation}

\textbf{关键物理洞察}：螺旋桨的平方负载（$Q \propto \omega^2$，其局部斜率为$2 k_Q \omega_0$）
本身就是天然的负反馈机制——当转速偏离工作点时，气动负载的变化倾向于将转速推回工作点。
这一机制在物理上等同于增大了等效阻尼系数，使$\tau_m$随$\omega_0$增大而减小（电机响应在高转速下更快）。

对于典型的小型无人机动力系统（$J_p \sim 2 \times 10^{-4}$\,kg$\cdot$m$^2$，
$k_t k_e / R \sim 2 \times 10^{-3}$\,N$\cdot$m$\cdot$s/rad），
螺旋桨负载的贡献$2 k_Q \omega_0$可与电磁阻尼同量级。
在高转速（高油门）下，$\tau_m$可降至低转速时的$60\%$--$70\%$——
这一效应在飞控增益调度中可通过油门-增益映射来近似补偿。

\subsection{旋转部件的陀螺效应}

旋转的电机-螺旋桨组件携带角动量$\bm{h} = J_p \omega \bm{n}$。
当机体有角速度$\bm{\Omega}$时，转子角动量矢量在惯性空间中的方向随体旋转而变化，
这一进动需要通过陀螺力矩来维持。由动系中的绝对导数公式（即矢量在不同参考系中时间导数的关系）：
\begin{equation}
    \left. \frac{d\bm{h}}{dt} \right|_{\text{惯性}}
    = \left. \frac{d\bm{h}}{dt} \right|_{\text{机体}}
    + \bm{\Omega} \times \bm{h}
    \label{eq:absolute_derivative}
\end{equation}

对于固连于机体的转子，在机体参考系中的角动量变化率$\dot{\bm{h}}|_{\text{机体}}$主要来自转速变化
（$J_p \dot{\omega} \bm{n}$）和摆座摆动导致的转轴方向变化。在大多数飞行动态中，
$\bm{\Omega} \times \bm{h}$项（进动效应）远大于$\dot{\bm{h}}|_{\text{机体}}$项
（转速变化的时间尺度通常比姿态变化慢）。忽略转子角动量大小的变化，得\textbf{陀螺力矩}：
\begin{equation}
    \boxed{\bm{M}_{\text{gyro}} = -\bm{\Omega} \times \bm{h}
           = -\bm{\Omega} \times (J_p \omega \bm{n})}
    \label{eq:gyro_moment}
\end{equation}

负号的含义是：陀螺力矩是机体角速度企图改变转子角动量方向时，转子对机体施加的反作用力矩——
本质上是对角动量方向改变的一种“抵抗”。

对于一个纵列双发构型，两台转子对机体的陀螺效应具有以下特征：
\begin{itemize}
    \item \textbf{机体俯仰角速率$q$}：前转子（转轴沿$+x_b$，$\bm{h}_f$沿$+x_b$）
          受$q$作用产生沿$+z_b$方向的进动趋势，陀螺力矩$-\bm{\Omega} \times \bm{h}_f$
          为$+z_b$方向力矩（偏航耦合）；尾转子（转轴沿$-x_b$，$\bm{h}_t$沿$-x_b$）
          的进动产生$-z_b$方向力矩——
          \textbf{前后转子的俯仰-偏航陀螺耦合恰好反向，理想情况下相互抵消}。
          这正是反向旋转设计的一个常被忽视的重要优势；
    \item \textbf{机体偏航角速率$r$}：前转子受$r$作用产生$-y_b$方向陀螺力矩（俯仰耦合），
          尾转子产生$+y_b$方向陀螺力矩——再次相互抵消；
    \item \textbf{机体滚转角速率$p$}：前转子转轴沿$x_b$，$\bm{\Omega} \times \bm{h}$为零
          （角速度矢量与角动量矢量平行），不产生陀螺力矩。尾转子同理。
    \item \textbf{摆座摆动}：摆座绕$z_b$（前）或$y_b$（尾）的摆动本身是转轴的附加角速度
          $\bm{\Omega}_{\text{gimbal}}$，在摆动瞬间产生附加的陀螺力矩
          $\bm{M} = -\bm{\Omega}_{\text{gimbal}} \times \bm{h}$，
          量级约为$J_p \omega \dot{\delta}$。除非摆动速率极高（如伺服的阶跃响应），
          该项通常可以忽略。
\end{itemize}

\textbf{结论}：在纵列双发、反向旋转的构型中，陀螺效应的一阶项在前后转子之间相互抵消。
这一抵消性质是选择反向旋转的重要物理理由，远超出仅为了静态反扭矩对消的直觉。


\section{推力矢量映射：从执行器到机体六维力/力矩}
